% New post-CP14 research fragment.  Not part of the frozen CP14 assembly.
\subsection{A superperfect-product obstruction to cohomological generation criteria}

\begin{example}[Proved: a superperfect-product obstruction]
\label{D:superperfect-quadratic-obstruction}
For every odd \(n\geq31\), let \(T_n\to A_n\) be the universal
central extension, and put
\[
 m_n=2^{n^2},\qquad
 \Pi=\prod_{\substack{n\geq31\\ n\ {\rm odd}}}T_n^{\,m_n}.
\]
Thus \(T_n\) is the double cover of \(A_n\).
The following assertions hold.
\begin{enumerate}
\item For every finite group \(H\), there are only finitely many continuous
homomorphisms \(\Pi\to H\).  Every open subgroup of \(\Pi\) has finite
maximal prosolvable quotient.
\item For every finite field \(k\), of any characteristic, and every
finite-dimensional continuous \(k[\Pi]\)-module \(M\), we have
\[
 \dim_k H^1(\Pi,M)\leq2\dim_k M,\qquad
 \dim_k H^2(\Pi,M)\leq5\dim_k M.
 \tag{SP1}
\]
For an open subgroup \(J\subset\Pi\), the corresponding constants
may be multiplied by \([\Pi:J]\).
\item There is one absolute constant \(C\), independent of \(q\), its
characteristic, and \(r\), such that
\[
 \#\left(
 \operatorname{Hom}_{\mathrm{cts}}(\Pi,\operatorname{GL}_r(\mathbf F_q))
 /\operatorname{GL}_r(\mathbf F_q)\right)\leq q^{Cr^2}.
 \tag{SP2}
\]
Nevertheless, \(\Pi\) is not topologically finitely generated.
\end{enumerate}
This is an abstract profinite example.  We assert no realization as
a geometric fundamental group.
\end{example}

\begin{proof}
We first specify the literature inputs and then prove the product
and counting assertions.  The Schur multiplier of \(A_n\) is \(C_2\)
for \(n\geq8\), and the universal central extension of a perfect group
has this multiplier as its kernel
\cite[Theorem~6.9.5 and Example~6.9.10, pp.~200--202]{Weibel1994SuperperfectB}.
We use this standard multiplier computation as a cited input.
The group \(T_n\) is perfect and has
\[
 H_1(T_n,\mathbf Z)=H_2(T_n,\mathbf Z)=0.
 \tag{SP3}
\]
For precision, the superperfectness assertion can be deduced from
universality as follows.  Let \(Y\to T_n\) be the universal central
extension of the perfect group \(T_n\); then \(Y\) is perfect.
For an element \(z\) in the kernel of \(Y\to A_n\), its commutators
with \(Y\) lie in the central kernel of \(Y\to T_n\).  Therefore
\(y\mapsto[z,y]\) is a homomorphism from the perfect group \(Y\)
to an abelian group and is zero.  Consequently \(Y\to A_n\) is
central.  The universality of \(T_n\to A_n\) splits \(Y\to T_n\);
the resulting central direct factor is trivial because \(Y\) is
perfect.  Thus \(Y=T_n\), proving (SP3).  This argument uses the
perfectness of \(Y\); a composite of arbitrary central extensions
need not be central.

The simplicity of \(A_n\) implies that every proper normal subgroup
of \(T_n\) lies in its central \(C_2\).  Indeed, if a normal subgroup
surjects onto \(A_n\), its product with the centre is \(T_n\), so
the corresponding quotient is abelian and hence trivial.
For odd \(n\), the elements
\((1\,2\,\cdots\,n)\) and \((1\,2\,3)\) generate \(A_n\):
their conjugates include the consecutive three-cycles, which
generate \(A_n\) by induction.  Any lifts of these two elements
generate \(T_n\).  Their generated subgroup \(L\) satisfies
\(LZ(T_n)=T_n\), hence is normal, and \(T_n/L\) is abelian;
perfectness gives \(L=T_n\).
In particular, for every \(T_n\)-module \(V\),
\[
 \dim H^1(T_n,V)\leq2\dim V,
 \tag{SP4}
\]
since a derivation is determined by its values on two generators.

\par\smallskip\noindent\textbf{The second-cohomology input.}\quad
We use the estimate
\[
 \dim H^2(A_n,V)<3\dim V
 \tag{SP5}
\]
from \cite[Theorem~6.2 and its proof in
Theorems~6.4, 6.5, 6.9, pp.~487--493]{GKKL2007SuperperfectB}.
The strict inequality can be weakened without affecting the argument.
The cited proof uses computational small-degree cases from its
reference~[42]; those computations are literature dependencies here,
not computations reproduced by our certificate.

We claim that \(\dim H^2(T_n,V)\leq5\dim V\) for all finite-dimensional
modules over every field.  In characteristic zero, positive-degree
cohomology of the finite group vanishes by averaging.
By coefficient extension and a composition
series, it suffices to consider an irreducible module over an
algebraically closed field \(k\).  The trivial module has \(H^2=0\)
by (SP3) and the universal-coefficient sequence.  Let \(V\) be
nontrivial.  If \(\operatorname{char}k\ne2\), averaging over the central
\(C_2\) gives either \(H^j(T_n,V)=0\) for every \(j\), when the centre
acts nontrivially, or
\(H^2(T_n,V)=H^2(A_n,V)\), when it acts trivially.
If \(\operatorname{char}k=2\), the normal central \(2\)-subgroup
acts trivially on \(V\).  The degree-two terms of the
Lyndon--Hochschild--Serre spectral sequence give
\[
 \dim H^2(T_n,V)
 \leq \dim H^2(A_n,V)+\dim H^1(A_n,V)+\dim V^{A_n}.
\]
The last term is zero, and \(A_n\) is two-generated, so (SP5)
gives the claimed bound.  This is also the central-subgroup
calculation in \cite[Lemma~4.1 and its proof, p.~481]
{GKKL2007SuperperfectB}.

\par\smallskip\noindent\textbf{Product cohomology and the infinite tail.}\quad
Over an algebraically closed field, every simple module for a finite
product of the \(T_n\)'s is a tensor product of simple modules of its
factors.  This holds in modular characteristic: quotienting the
finite-dimensional group algebras by their radicals reduces the
claim to tensor products of matrix algebras.
Tensoring free bar resolutions and applying the Kunneth formula
over a field gives
\[
 H^a\!\left(\prod_i G_i,\bigotimes_i V_i\right)
 =
 \bigoplus_{\sum_i a_i=a}\bigotimes_i H^{a_i}(G_i,V_i).
 \tag{SP6}
\]
The trivial coefficient has \(H^1(T_n,k)=H^2(T_n,k)=0\).
For a nontrivial simple factor \(V_i\), \(H^0(T_n,V_i)=0\).
It follows that a simple module with one nontrivial support factor
has \(H^1\) of dimension at most twice its dimension and \(H^2\)
of dimension at most five times its dimension.  With two nontrivial
support factors, \(H^1=0\), and the only possible degree-two term is
the tensor product of two \(H^1\)'s, of dimension at most four times
the module dimension.  With at least three support factors,
both \(H^1\) and \(H^2\) vanish.  For the trivial simple module both
groups vanish.

These statements compute the continuous cohomology of the full
infinite product as well.  A finite coefficient module is trivial
on a full tail.  Continuous cochains with finite values on that
tail factor through a finite coordinate product, so its cochain
complex is the filtered union of the corresponding finite-product
complexes.  Formula (SP6) and (SP3) show that the tail has zero
\(H^1\) and \(H^2\) with trivial coefficients.  The
Lyndon--Hochschild--Serre spectral sequence for head times tail
therefore identifies \(H^1\) and \(H^2\) with those of the finite
head acting on the given module.  Coefficient extension and
devissage now prove (SP1).  Shapiro's lemma proves the stated
open-subgroup bounds, since induction from \(J\) multiplies the
coefficient dimension by \([\Pi:J]\).

\par\smallskip\noindent\textbf{A characteristic-independent dimension bound.}\quad
Let \(V\) be a nontrivial simple \(T_n\)-module over an algebraically
closed field of characteristic \(\ell\).  Its kernel is contained
in the centre.  If \(\ell\ne3\), take the elementary abelian subgroup
of \(A_n\) generated by \(\lfloor n/3\rfloor\) disjoint three-cycles.
If \(\ell=3\), take \(\lfloor n/5\rfloor\) disjoint five-cycles.
Each of these odd-order elementary abelian groups lifts to a
subgroup of \(T_n\) with the same order.  To check this directly,
let its exponent be the odd prime \(s\).  A lift of a basis
element can be multiplied by the central involution to have order
\(s\).  Commutators of these lifts are central of order dividing
two and also of order dividing \(s\), so they commute.
The resulting elementary abelian subgroup meets the centre
trivially and acts faithfully on \(V\).

Its order is prime to \(\ell\), hence its action is simultaneously
diagonalizable.  A faithful diagonal representation of
\((C_s)^a\) requires at least \(a\) characters, since they must
span its character space.  Consequently
\[
 \dim V\geq\lfloor n/5\rfloor\geq n/6,\qquad n\geq31.
 \tag{SP7}
\]
In particular every nontrivial simple factor has dimension at least
two.  This argument avoids using a lifted elementary abelian
\(2\)-subgroup, which need not remain abelian.

\par\smallskip\noindent\textbf{Counting simple modules.}\quad
Let an absolutely simple module of dimension \(r\geq2\) have
nontrivial support dimensions \(a_1,\ldots,a_t\), at degrees
\(n_1,\ldots,n_t\).  Then
\[
 r=\prod_i a_i,\qquad t\leq\log_2r,\qquad
 \sum_i a_i^2\leq r^2,\qquad \sum_i n_i^2\leq36r^2.
 \tag{SP8}
\]
The middle inequality follows by induction from
\(a^2+b^2\leq a^2b^2\) for \(a,b\geq2\).
We may order the support, thereby overcounting.  There are at most
\((6r)^{1+\lfloor\log_2r\rfloor}\leq2^{2r^2}\)
choices for its degree sequence.
Here the displayed estimate follows from
\(1+\lfloor\log_2r\rfloor\leq r\) and
\(\log_2(6r)\leq2r\) for \(r\geq2\).
The coordinate-copy choices
number at most \(2^{\sum_i n_i^2}\).
For each coordinate, the number of simple labels is at most
\(\dim k[T_{n_i}]=|T_{n_i}|=n_i!\leq2^{n_i^2}\).
Thus the number of absolutely simple \(r\)-dimensional modules is
at most \(2^{74r^2}\), uniformly in the coefficient characteristic.
For \(r=1\) there is only the trivial module.

For a finite field \(\mathbf F_q\), every simple module becomes
over its algebraic closure a sum of distinct Frobenius-conjugate
absolutely simple modules, with unique descent.  Indeed, the
semisimple quotient of a finite group algebra is a product of
matrix algebras over finite fields, since finite division rings
are commutative.  Splitting each field factor gives precisely
this description.  Hence the number \(s_r(q)\) of simple
\(\mathbf F_q[\Pi]\)-modules of dimension \(r\) satisfies
\[
 s_r(q)\leq\sum_{e\mid r}2^{74(r/e)^2}\leq q^{75r^2}.
 \tag{SP9}
\]
The first inequality may overcount but is enough.

\par\smallskip\noindent\textbf{Extensions and arbitrary representations.}\quad
For two absolutely simple modules \(U,V\), tensoring bar resolutions
gives
\[
 \operatorname{Ext}^1(U,V)=
 \bigoplus_i H^1(G_i,\operatorname{Hom}(U_i,V_i))
 \otimes\!\!\bigotimes_{j\ne i}\operatorname{Hom}_{G_j}(U_j,V_j).
 \tag{SP10}
\]
If they differ at two coordinates, this is zero.  If they differ
at one coordinate, (SP4) bounds its dimension by
\(2\dim U\dim V\).  If \(U=V\), Schur's lemma and (SP8)
bound it by \(2\sum_i a_i^2\leq2(\dim U)^2\).
Trivial coordinates contribute no term because the coordinate
groups are perfect.  In particular, the computation remains valid
for the full product: an extension cannot acquire new support
on a factor trivial on both its submodule and quotient, since
such a factor would act through an additive abelian group.
After coefficient extension and devissage in the second variable,
\[
 \dim_{\mathbf F_q}\operatorname{Ext}^1(U,W)
 \leq2\dim U\dim W
 \tag{SP11}
\]
for every simple \(U\) and every finite-dimensional \(W\).

Fix a composition \(r=r_1+\cdots+r_t\).  By (SP9) there are at
most \(q^{75\sum_i r_i^2}\) choices of simple factors.  Successive
extensions contribute at most \(q^{2\sum_{i<j}r_ir_j}\), by (SP11).
We count extensions with identified ends, which only overcounts
isomorphism classes.  Since
\[
 75\sum_i r_i^2+2\sum_{i<j}r_ir_j\leq75r^2
\]
and there are \(2^{r-1}\) compositions of \(r\), we obtain (SP2),
for example with \(C=76\).  Isomorphism classes of these modules
are exactly conjugacy classes of the representations in (SP2).

\par\smallskip\noindent\textbf{Smallness and failure of finite generation.}\quad
A nontrivial image of \(T_n\) has order at least \(|A_n|\).
Thus a homomorphism to a fixed finite group kills all but finitely
many degrees, and only finitely many coordinate groups remain.
Continuity kills their full tail, proving the asserted finiteness
of homomorphisms.  An open subgroup contains a full tail and is
a product of that tail with a subgroup of a finite head.
Every finite solvable quotient kills the perfect tail, so the
maximal prosolvable quotient is finite.

Finally, \(\Pi\) maps onto \(A_n^{m_n}\) for every odd \(n\geq31\).
If \(d\) elements generated \(\Pi\), the coordinate \(d\)-tuples
in this quotient would have to be pairwise distinct: equal
tuples force a diagonal relation at the corresponding two
coordinates.  Therefore \(m_n\leq|A_n|^d\).
This contradicts \(2^{n^2}>(n!)^d\) for all sufficiently large
odd \(n\), proving non-finite-generation.
\end{proof}

\begin{remark}[Literature and exact scope]
Products of Schur covers of alternating groups already occur in
\cite[Section~1.6, Proposition~1.14, and Section~2, especially
Proposition~2.4 and Theorem~2.5]{CorobCookVannacci2021B}.
We use this literature as context and make no claim of priority
for the construction or its variants.  The present multiplicities
are \(2^{n^2}\), and the assertions proved above are the precise
ones used here.
The example shows that adding an all-characteristic linear
\(H^2\) bound to smallness, finite open prosolvable quotients,
linear \(H^1\), and uniform quadratic representation growth
still does not by itself imply topological finite generation.
It does not refute any geometric statement in this paper.
The finite exact certificate
\nolinkurl{parts/B/post14_research/code/check_superperfect_product.py}
checks the displayed permutation generators, odd-cycle subgroups,
support-square and counting inequalities, and the low-degree tensor
coefficient bounds.  Its output is
\nolinkurl{parts/B/post14_research/output/superperfect_product.json}.
It does not compute Schur multipliers or repeat the cohomology
computations cited in the literature.
\end{remark}
